\section{Data dictionary}\label{sec:diccionario}

This appendix documents the consolidated dataset produced by
\texttt{src/importar\_datos.py} from the six annual SIMA workbooks
(\texttt{data/raw/BD 2020.xlsx} to \texttt{data/raw/BD 2025.xlsx}): what each
variable contains, how the stations are covered, and which quality anomalies
must be addressed before any multivariate analysis.

\subsection{Size of the dataset}\label{sec:dic:dimension}

The consolidated dataset has 754\,603 rows and 18 columns: 3 identifiers and
15 measured parameters. Each row is one station in one hour, coverage runs
from 1 January 2020 at 00:00 to 31 December 2025 at 23:00, and there are no
duplicate records. Of the 11\,319\,045 parameter cells, 1\,188\,343 are
empty: 10.50\,\% missing.

At the source, each annual workbook has one sheet per station ---13 in 2020,
14 in 2021 and 15 from 2022 onward--- in wide format: one row per hour and one
column per parameter. There is no station column; consolidation promotes the
sheet name to the \texttt{estacion} column and stacks the six years.

\subsection{Identifier variables}\label{sec:dic:identificadoras}

\begin{description}
  \item[\texttt{fecha}] Timestamp of the observation (date), at hourly
    resolution (\texttt{datetime64}), from 2020-01-01 00:00 to 2025-12-31
    23:00. No missing values.
  \item[\texttt{estacion}] Nominal categorical variable (station) with 15
    levels: CE, NE, NE2, NE3, NO, NO2, NO3, NTE, NTE2, SE, SE2, SE3, SO, SO2
    and SUR. No missing values.
  \item[\texttt{anio}] Calendar year, derived from the source file. Discrete
    numeric (or ordinal categorical), with values from 2020 to 2025. No
    missing values.
\end{description}

\subsection{Criteria pollutants}\label{sec:dic:contaminantes}

The eight criteria pollutants are continuous numeric variables and are
summarized in \cref{tab:dic-contaminantes}. The reported operating range spans
from the most restrictive to the most permissive year of the period, because
the limit changes from year to year (\cref{sec:dic:rangos}). According to
\texttt{Etiquetas.xlsx}, O$_3$, SO$_2$ and NO$_2$ can be converted to ppm by
dividing by 1000.

\subsection{Meteorological parameters}\label{sec:dic:meteorologicos}

The seven meteorological parameters appear in \cref{tab:dic-meteo}. Two of
them cannot be treated as ordinary continuous variables:

\begin{itemize}
  \item \texttt{WDR} is a \textbf{circular} variable, even though it is stored
    as a number: $0^{\circ}$ and $360^{\circ}$ are the same direction (north).
    It cannot be averaged, scaled or correlated like the others; it must be
    decomposed into $\sin\theta$ and $\cos\theta$, or into the wind components
    $u = -\,\mathrm{WSR}\cdot\sin\theta$ and
    $v = -\,\mathrm{WSR}\cdot\cos\theta$. It is a derived attribute that must
    be built at the transformation stage.
  \item \texttt{RAINF} and \texttt{SR} are \textbf{zero-inflated}: 98.4\,\% of
    the valid precipitation records and 39.6\,\% of the solar radiation
    records are exactly zero (night and days without rain). These are
    legitimate zeros, not missing values or errors, but they break normality
    assumptions and distort any standardization.
\end{itemize}

\subsection{Station catalog}\label{sec:dic:estaciones}

\Cref{tab:dic-estaciones} lists the 15 stations with the description from
\texttt{Etiquetas.xlsx} and the coverage computed on the consolidated
dataset. Two stations are not documented in the SIMA catalog and also joined
the network late: NE3 appears from 2021 and NO3 only from December 2022, so
any analysis that requires the full panel must exclude them or be trimmed to
the common period.

\subsection{Quality flags}\label{sec:dic:banderas}

SIMA documents in \texttt{Etiquetas.xlsx} a system of flags that records why
an hour is invalidated (\cref{tab:dic-banderas}). These flags \textbf{are not
included in the \texttt{BD *.xlsx} files}: the delivered workbooks contain
only the numeric value or an empty cell. The data arrive already filtered;
each empty cell is an hour that SIMA invalidated for one of those reasons, but
the reason cannot be recovered.

The methodological consequence is direct: the missing values in this dataset
\textbf{are not random} (not MCAR). They concentrate where there was
calibration, failure or an anomalous condition, which brings them closer to a
MAR mechanism. This constrains which imputation is defensible, and it must be
declared as an assumption.

\subsection{Missing data}\label{sec:dic:faltantes}

\Cref{tab:dic-faltantes} gives the percentage of missing values per variable,
which ranges from 3.66\,\% for \texttt{SR} to 25.23\,\% for \texttt{PM2.5}.
That average hides what matters: missingness is not evenly distributed across
stations or years.

For ozone, \textbf{2020 is practically unusable}
(\cref{tab:dic-o3-faltante}): NTE, SUR and SE2 have 100\,\% of ozone missing
that year, NE 99.9\,\%, NO 98.4\,\%, NTE2 93.3\,\% and NE2 89.1\,\%. From
2022 onward, missing O$_3$ falls to between 1\,\% and 8\,\% across almost the
entire network; the exception is NO3, with 30.3\,\% in 2025.

\subsection{Sentinel and physically impossible values}\label{sec:dic:centinela}

\Cref{tab:dic-anomalias} gathers the values that cannot be taken as
measurements. The clearest case is $-9999$, the missing-data code of the
\emph{datalogger}: 26 cells spread across NOX, RAINF, RH, SR, TOUT, WSR and
WDR are enough to shift means and correlations, so they must be converted to
\texttt{NaN} before any computation.

\subsection{Operating ranges are not stable across years}\label{sec:dic:rangos}

The SIMA range document defines different limits for each year, and some are
stricter than the actual data. The critical case is \texttt{PRS}: the range
declared for 2022 is 700--740 mmHg, and with that criterion 6\,305 records of
that year (about 4.8\,\%) would fall outside it, whereas in 2021, 2023, 2024
and 2025 the same kind of reading is valid because the lower limit drops to
687--690. Applying the document literally would delete legitimate data.

The same happens with \texttt{TOUT}, whose declared minimum is $0^{\circ}$C
in 2020 and 2023 but $-6.5^{\circ}$C in 2021: Monterrey has real frosts, so
the $0^{\circ}$C limit is a data-entry error in the document and not a
physical rule.

The adopted criterion is to use the \textbf{manufacturer range}
(\cref{tab:dic-fabricante}), which is physical and stable, as a hard filter
to remove impossible values, and to treat the annual operating range as a
signal of suspicion to review, not as an automatic deletion rule.

\subsection{Incomplete time grid}\label{sec:dic:rejilla}

No station has a complete hourly series: between 3 (CE) and 16 (SUR)
timestamps are missing relative to the continuous hourly grid of its period.
These are missing hours, not empty rows: the row simply does not exist. For
this reason, any temporal interpolation requires reindexing against a
complete hourly grid first; otherwise the method will assume that two
contiguous observations are one hour apart when they are not.

\begin{table*}[t]
  \centering
  \caption{Criteria pollutants. All are continuous numeric variables. The
    operating range goes from the most restrictive to the most permissive
    year of the period.}
  \label{tab:dic-contaminantes}
  \small
  \begin{tabular}{lllcccc}
    \toprule
    Variable & Description & Unit & Operating range & Observed range & Median & \% missing \\
    \midrule
    \texttt{CO}    & Carbon monoxide                       & ppm             & 0--8 to 0--20       & $0.0$--$37.0$    & 1.17  & 13.31 \\
    \texttt{NO}    & Nitric oxide                          & ppb             & 0--350 to 0--500    & $0.3$--$945.1$   & 5.00  & 16.75 \\
    \texttt{NO2}   & Nitrogen dioxide                      & ppb             & 0--100 to 0--200    & $0.0$--$188.6$   & 11.40 & 17.37 \\
    \texttt{NOX}   & Nitrogen oxides (NO $+$ NO$_2$)       & ppb             & 0--400 to 0--500    & $-9999$--$971.8$ & 17.30 & 16.82 \\
    \texttt{O3}    & \textbf{Ozone}: target variable       & ppb             & 0--153 to 0--185    & $0.5$--$265.0$   & 23.00 & 14.73 \\
    \texttt{PM10}  & Particles $<10\,\mu$m                 & $\mu$g/m$^3$    & 0--800 to 0--999    & $1.0$--$1001.0$  & 49.00 & 4.63  \\
    \texttt{PM2.5} & Particles $<2.5\,\mu$m                & $\mu$g/m$^3$    & 0--205.94 to 0--999 & $0.0$--$999.0$   & 17.00 & 25.23 \\
    \texttt{SO2}   & Sulfur dioxide                        & ppb             & 0--150 to 0--405    & $0.0$--$404.7$   & 3.80  & 15.03 \\
    \bottomrule
  \end{tabular}
\end{table*}

\begin{table*}[t]
  \centering
  \caption{Meteorological parameters. All are continuous numeric variables,
    with the caveats of \texttt{WDR} (circular) and of \texttt{RAINF} and
    \texttt{SR} (zero-inflated).}
  \label{tab:dic-meteo}
  \small
  \begin{tabular}{lllcccc}
    \toprule
    Variable & Description & Unit & Operating range & Observed range & Median & \% missing \\
    \midrule
    \texttt{TOUT}  & Ambient temperature  & $^{\circ}$C        & $-6.5$--45 to 0--45.5 & $-9999$--$785.0$ & 24.02  & 6.95  \\
    \texttt{RH}    & Relative humidity    & \%                 & 0--100                & $-9999$--$725.0$ & 59.00  & 11.38 \\
    \texttt{SR}    & Solar radiation      & kW/m$^2$           & 0--1 to 0--1.26       & $-9999$--$604.0$ & 0.00   & 3.66  \\
    \texttt{RAINF} & Precipitation        & mm/h               & 0--25 to 0--80        & $-9999$--$360.0$ & 0.00   & 4.74  \\
    \texttt{PRS}   & Atmospheric pressure & mmHg               & 687.5--750            & $0.0$--$750.0$   & 714.40 & 5.08  \\
    \texttt{WSR}   & Wind speed           & km/h               & 0--35 to 0--75        & $-9999$--$337.0$ & 7.40   & 6.87  \\
    \texttt{WDR}   & Wind direction       & azimuth degrees    & 0--360                & $-9999$--$360.0$ & 116.00 & 8.35  \\
    \bottomrule
  \end{tabular}
\end{table*}

\begin{table*}[t]
  \centering
  \caption{Catalog of monitoring stations. The description comes from
    \texttt{Etiquetas.xlsx}; coverage and hours were computed on the
    consolidated dataset.}
  \label{tab:dic-estaciones}
  \small
  \begin{tabular}{lllllr}
    \toprule
    Code & Name & Site & Municipality & Coverage & Hours \\
    \midrule
    \texttt{CE}   & Centro                & Obispado       & Monterrey                & 2020--2025                & 52\,605 \\
    \texttt{NE}   & Noreste               & San Nicolás    & San Nicolás de los Garza & 2020--2025                & 52\,595 \\
    \texttt{NE2}  & Noreste 2             & Apodaca        & Apodaca                  & 2020--2025                & 52\,593 \\
    \texttt{NE3}  & \emph{Undocumented}   & ---            & ---                      & \textbf{2021--2025}       & 43\,817 \\
    \texttt{NO}   & Noroeste              & San Bernabé    & Monterrey                & 2020--2025                & 52\,597 \\
    \texttt{NO2}  & Noroeste 2            & García         & García                   & 2020--2025                & 52\,594 \\
    \texttt{NO3}  & \emph{Undocumented}   & ---            & ---                      & \textbf{Dec.\ 2022--2025} & 27\,045 \\
    \texttt{NTE}  & Norte                 & Escobedo       & Escobedo                 & 2020--2025                & 52\,595 \\
    \texttt{NTE2} & Norte 2               & Universidad    & (not given)              & 2020--2025                & 52\,593 \\
    \texttt{SE}   & Sureste               & La Pastora     & Guadalupe                & 2020--2025                & 52\,599 \\
    \texttt{SE2}  & Sureste 2             & Juárez         & Juárez                   & 2020--2025                & 52\,595 \\
    \texttt{SE3}  & Sureste 3             & Cadereyta      & (not given)              & 2020--2025                & 52\,593 \\
    \texttt{SO}   & Suroeste              & Santa Catarina & Santa Catarina           & 2020--2025                & 52\,597 \\
    \texttt{SO2}  & Suroeste 2            & San Pedro      & San Pedro Garza García   & 2020--2025                & 52\,593 \\
    \texttt{SUR}  & Sur                   & Pueblo Serena  & (not given)              & 2020--2025                & 52\,592 \\
    \bottomrule
  \end{tabular}
\end{table*}

\begin{table*}[t]
  \centering
  \caption{Quality flags documented in \texttt{Etiquetas.xlsx}. They do not
    appear in the \texttt{BD *.xlsx} workbooks, so the reason for
    invalidation cannot be recovered from the delivered data.}
  \label{tab:dic-banderas}
  \small
  \begin{tabular}{llcc}
    \toprule
    Flag & Cause & Valid-hour mark & Invalid-hour mark \\
    \midrule
    \texttt{P} / \texttt{p} & Power failure                                      & \texttt{P} & \texttt{p} \\
    \texttt{C} / \texttt{c} & Calibration                                        & \texttt{C} & \texttt{c} \\
    \texttt{D} / \texttt{d} & Shutdown                                           & \texttt{D} & \texttt{d} \\
    \texttt{B} / \texttt{b} & Bad conditions                                     & \texttt{B} & \texttt{b} \\
    \texttt{m}              & Positive above range                               & ---        & \texttt{m} \\
    \texttt{l}              & Negative below range or solar radiation adjustment & ---        & \texttt{l} \\
    \texttt{z}              & Zeros and negatives                                & ---        & \texttt{z} \\
    \texttt{o}              & PM10 above $900\,\mu$g/m$^3$                       & ---        & \texttt{o} \\
    \texttt{s}              & Repeated or identical consecutive values           & ---        & \texttt{s} \\
    \texttt{r}              & PM10 versus PM2.5 comparison                       & ---        & \texttt{r} \\
    \texttt{e}              & Delete NO and NOx data                             & ---        & \texttt{e} \\
    \texttt{a}              & Delete PM $<5\,\mu$g/m$^3$ and CO $<0.05$ ppm      & ---        & \texttt{a} \\
    \texttt{f}              & Value $3\times$ the previous one (PM10)            & ---        & \texttt{f} \\
    \texttt{h}              & Jump $>10^{\circ}$C or $>10$ mmHg from the previous hour & --- & \texttt{h} \\
    \texttt{n}              & Communication failure                              & ---        & \texttt{n} \\
    \texttt{x}              & Value that passed the protocols but should not be present & --- & \texttt{x} \\
    \texttt{u}              & Undocumented                                       & \multicolumn{2}{c}{Undocumented} \\
    \bottomrule
  \end{tabular}
\end{table*}

\begin{table}[t]
  \centering
  \caption{Percentage of missing values per variable.}
  \label{tab:dic-faltantes}
  \small
  \begin{tabular}{lr@{\qquad}lr}
    \toprule
    Variable & \% missing & Variable & \% missing \\
    \midrule
    \texttt{PM2.5} & 25.23 & \texttt{RH}    & 11.38 \\
    \texttt{NO2}   & 17.37 & \texttt{WDR}   & 8.35  \\
    \texttt{NOX}   & 16.82 & \texttt{TOUT}  & 6.95  \\
    \texttt{NO}    & 16.75 & \texttt{WSR}   & 6.87  \\
    \texttt{SO2}   & 15.03 & \texttt{PRS}   & 5.08  \\
    \texttt{O3}    & 14.73 & \texttt{RAINF} & 4.74  \\
    \texttt{CO}    & 13.31 & \texttt{PM10}  & 4.63  \\
                   &       & \texttt{SR}    & 3.66  \\
    \bottomrule
  \end{tabular}
\end{table}

\begin{table}[t]
  \centering
  \caption{Percentage of missing O$_3$ by station and year. A dash indicates
    that the station was not yet operating.}
  \label{tab:dic-o3-faltante}
  \small
  \begin{tabular}{lrrrrrr}
    \toprule
    Station & 2020 & 2021 & 2022 & 2023 & 2024 & 2025 \\
    \midrule
    \texttt{CE}   & 16.5  & 12.2 & 3.6 & 4.1  & 4.2  & 3.4  \\
    \texttt{NE}   & 99.9  & 16.0 & 3.8 & 5.0  & 3.7  & 10.1 \\
    \texttt{NE2}  & 89.1  & 18.6 & 3.3 & 7.8  & 7.0  & 6.0  \\
    \texttt{NE3}  & ---   & 4.9  & 4.4 & 4.2  & 5.2  & 7.6  \\
    \texttt{NO}   & 98.4  & 20.3 & 6.2 & 6.0  & 11.0 & 7.4  \\
    \texttt{NO2}  & 20.3  & 5.9  & 2.4 & 2.7  & 3.8  & 4.7  \\
    \texttt{NO3}  & ---   & ---  & 5.1 & 16.5 & 3.5  & 30.3 \\
    \texttt{NTE}  & 100.0 & 18.5 & 4.3 & 7.3  & 4.0  & 4.2  \\
    \texttt{NTE2} & 93.3  & 9.5  & 2.5 & 2.9  & 5.9  & 4.6  \\
    \texttt{SE}   & 7.3   & 5.6  & 2.7 & 2.8  & 2.8  & 3.0  \\
    \texttt{SE2}  & 100.0 & 39.3 & 5.5 & 1.6  & 12.0 & 10.0 \\
    \texttt{SE3}  & 28.5  & 5.9  & 1.8 & 4.7  & 1.4  & 0.5  \\
    \texttt{SO}   & 31.0  & 7.0  & 5.4 & 4.3  & 5.9  & 4.7  \\
    \texttt{SO2}  & 3.0   & 2.5  & 2.6 & 1.9  & 1.4  & 1.4  \\
    \texttt{SUR}  & 100.0 & 16.4 & 2.6 & 2.8  & 0.9  & 0.9  \\
    \bottomrule
  \end{tabular}
\end{table}

\begin{table*}[t]
  \centering
  \caption{Sentinel and physically impossible values detected in the
    consolidated dataset.}
  \label{tab:dic-anomalias}
  \small
  \begin{tabular}{p{3cm}p{2.6cm}rp{7cm}}
    \toprule
    Anomaly & Affected variables & $n$ & Interpretation \\
    \midrule
    $-9999$
      & NOX, RAINF, RH, SR, TOUT, WSR, WDR & 26
      & Missing-data code of the \emph{datalogger}, \textbf{not a value}. Must be set to \texttt{NaN} before any computation. \\
    \addlinespace
    $\mathrm{RH} > 100$\,\% & RH & 285
      & Impossible by definition; maxima of 714\,\% and 725\,\% are observed. \\
    \addlinespace
    $\mathrm{TOUT} > 50^{\circ}$C & TOUT & 3
      & Outside the sensor range ($-50$ to $50^{\circ}$C); the observed maximum is $785^{\circ}$C. \\
    \addlinespace
    \texttt{SR} on the wrong scale & SR & $\approx 30$
      & Values of 64, 501 and 604 against a physical maximum of 1.4 kW/m$^2$: readings in W/m$^2$ mixed in (divide by 1000). \\
    \addlinespace
    \texttt{SR} between 1.5 and 7.8 & SR & $\approx 730$
      & Exceed the manufacturer range and concentrate at NTE (493) and SO (172), which suggests sensor failure or miscalibration rather than a meteorological event. \\
    \addlinespace
    $\mathrm{WSR} > 100$ km/h & WSR & 418
      & Above any wind recorded in the MMA; 306 of them occur in 2021. \\
    \addlinespace
    $\mathrm{PRS} = 0$ or $\mathrm{PRS} < 500$ & PRS & 2
      & A zero atmospheric pressure is impossible. \\
    \addlinespace
    $\mathrm{PM10} \geq 999$, $\mathrm{PM2.5} = 999$ & PM10, PM2.5 & 6
      & Cap values (999, 1000, 1001) characteristic of saturation or sentinel codes, not real concentrations. \\
    \addlinespace
    $\mathrm{CO} > 18$ ppm & CO & 24
      & Exceeds the operating range of every year; the maximum is 37 ppm. \\
    \bottomrule
  \end{tabular}
\end{table*}

\begin{table}[t]
  \centering
  \caption{Manufacturer range per parameter, adopted as a hard filter for
    impossible values.}
  \label{tab:dic-fabricante}
  \small
  \begin{tabular}{lc}
    \toprule
    Parameter & Manufacturer range \\
    \midrule
    PM10, PM2.5, O3      & $0$--$1000$     \\
    NO, NO2, NOx, SO2    & $0$--$500$      \\
    CO                   & $0$--$50$       \\
    RH                   & $0$--$100$      \\
    WS (\texttt{WSR})    & $0$--$180$      \\
    TEMP (\texttt{TOUT}) & $-50$--$50$     \\
    SR                   & $0$--$1.4$      \\
    BP (\texttt{PRS})    & $449.9$--$824.9$\\
    WD (\texttt{WDR})    & $0$--$360$      \\
    \bottomrule
  \end{tabular}
\end{table}
